\section{AutoResearch：递归自我改进的试验田}

\subsection{核心理念：把自己从循环中移除}

Karpathy 之前发过一条推文，成为 AutoResearch 的哲学基础：

\begin{quote}
``To get the most out of the tools that have become available now you have to remove yourself as the bottleneck.''

——要充分利用已有工具的全部潜力，你必须把自己作为瓶颈移除。
\end{quote}

你不能在那里等着提示下一步。你要安排好一切，让系统完全自主运行。当下竞争的本质是增加你的杠杆——你偶尔投入一点 Token，大量的事情代替你发生。

\subsection{nanochat：两个十年的经验 vs 一个晚上的 Agent}

很多人对 Karpathy 痴迷训练 GPT-2 级别的模型感到困惑，但对他来说，这是\textbf{递归自我改进（recursive self-improvement）}的试验田——正是所有前沿实验室都在追求的方向。

他已经用"老派"方式把 nanochat（他的精简 LLM 训练框架）调得相当好——做了二十年研究，做了大量超参数搜索和实验。然后他让 AutoResearch 跑了一个晚上。

Agent 回来的时候带着他没发现的优化：
\begin{itemize}
    \item Value embeddings 上的 weight decay 漏了
    \item Adam 优化器的 beta 参数没调够
    \item 这些参数之间有联合交互效应——调了一个，其他的最优值也变了
\end{itemize}

\begin{warningbox}{过度自信的陷阱}
Karpathy 坦言自己有"两个十年的earned confidence"，自认为已经调好了。但 AutoResearch 一个晚上就发现了他遗漏的优化。他的结论是：研究者不应该成为瓶颈——不应该是你在跑超参数搜索、不应该是你在看结果。有客观标准的地方，就应该让系统自主运行。
\end{warningbox}

\begin{knowledgebox}{AutoResearch 公开成果}
AutoResearch 于2026年3月7日开源，核心是一个约630行的 Python 脚本，让 AI Agent 在单 GPU 上自主循环运行实验。两天连续运行约700次实验，发现约20个叠加有效的改进，将"Time to GPT-2"排行榜指标从2.02小时降到1.80小时——在一个已经调好的项目上获得了11\%的效率提升。Shopify CEO Tobias L\"{u}tke 用同样的方法在自己公司内部数据上跑了37个实验，获得了19\%的性能提升。
\end{knowledgebox}

\subsection{program.md：每个研究组织都是一组 Markdown 文件}

Sarah 问：什么时候模型能写出比你更好的 program.md？

Karpathy 展开了一个框架：

\begin{quote}
``Every research organization is described by program.md. A research organization is a set of markdown files.''

——每个研究组织都是由 program.md 描述的。一个研究组织就是一组 Markdown 文件。
\end{quote}

不同的 program.md 会产生不同的研究进展。一个组织可以少开晨会（因为没用），另一个可以多承担风险。你可以想象让多个"研究组织"竞赛，然后分析改进来自哪里，用分析结果让模型生成更好的 program.md——这就是\textbf{元优化（meta-optimization）}。

Sarah 提出了一个竞赛想法：让不同的人写不同的 program.md，在相同硬件上比谁获得最大改进，然后把所有数据交给模型写出更好的 program.md。Karpathy 认为这100\%可行——"我们一定会得到更好的东西。"

\subsection{从单循环到分布式：Open Ground 的设想}

AutoResearch 当前是单线程的——一个 Agent 在循环中不断改进。Karpathy 一直在思考如何并行化，特别是引入互联网上的\textbf{不可信工作者（untrusted workers）}。

他的设计思路类似区块链：
\begin{itemize}
    \item \textbf{Commits} 对应区块——任何人都可以从互联网上提交一个 commit，声称这段代码能优化性能
    \item \textbf{实验} 对应工作量证明——产生一个好 commit 需要大量搜索工作（可能尝试了一万个想法，只有一个成功）
    \item \textbf{验证便宜但搜索昂贵}——你训练一次就知道 commit 是否有效，但产生它可能需要上千次尝试
\end{itemize}

\begin{quote}
``A swarm of agents on the internet could collaborate to improve LLMs and could potentially even run circles around frontier labs.''

——互联网上的 Agent 集群有可能协作改进 LLM，甚至可能绕着前沿实验室跑圈。
\end{quote}

逻辑是：前沿实验室有大量可信算力，但地球更大，有海量不可信算力。如果你能设计好验证系统，分布式集群未必不能胜出。

\begin{knowledgebox}{类比：SETI@Home 与 Folding@Home}
SETI@Home 和 Folding@Home 是两个著名的分布式计算项目，前者搜索外星文明信号，后者模拟蛋白质折叠。都利用全球志愿者的闲置电脑算力。AutoResearch 的分布式版本遵循相同的架构：产生结果很难，但验证结果很便宜。
\end{knowledgebox}

他还进一步设想：如果一切都被重构为 AutoResearch，算力就是你贡献给公共事业的新货币。你关心某种类型的癌症研究？不用只是捐钱——你可以购买算力，加入那个项目的 AutoResearch 池。

\subsection{本章小结}

AutoResearch 的核心是"把人从循环中移除"。单循环版本已经在调好的项目上发现了11\%的优化空间。program.md 框架将研究组织抽象为可优化的代码。分布式版本（Open Ground）如果成功，可能根本性地改变前沿研究的格局——从中心化实验室转向互联网 Agent 集群。


